% !TeX spellcheck = ru_RU
% !TEX root = vkr.tex

\section*{Введение}
\thispagestyle{withCompileDate}
Современные языки программирования представляют собой не только синтаксис и набор семантических правил, но и программную инфраструктуру, обеспечивающую выполнение написанных на них программ~\cite{AhoCompilers, EngineeringCompiler}. В такую инфраструктуру обычно входят компилятор, промежуточные представления программы, среда исполнения и один или несколько механизмов запуска программы. От полноты и согласованности этих компонентов зависит не только удобство использования языка, но и возможность его дальнейшего развития.

Одним из распространенных подходов  к организации исполнения программ является использование байткода и виртуальной машины~\cite{JVMspec, PythonVM, LuaVM}. В такой архитектуре исходная программа сначала преобразуется в промежуточное низкоуровневое представление, не привязанное напрямую к конкретному процессору, а затем выполняется специальной средой исполнения. Этот подход используется как в промышленных языках программирования, так и в учебных и исследовательских проектах, поскольку он позволяет отделить компиляцию от конкретного способа выполнения программы.

Lama является учебным языком программирования, разрабатываемый для знакомства с областью языков
программирования, компиляторов и инструментов~\cite{LamaReadme,LamaSpec}. Язык сочетает императивные
и функциональные конструкции: переменные, функции, рекурсию, замыкания, массивы, строки, S-выражения
и сопоставление с образцом. В существующей реализации Lama исходная программа компилируется в
промежуточное стековое представление, после чего может быть либо скомпилирована в нативный код, либо
в байткод.

% TODO: про переносимость и почему существующий интепретатор плох и почему нужен Си
Тем не менее, на данный момент не существует стандартного интерпретатора для байткода.
Кроме того, формат байткода неспецифированный и его запуск многомодульных программ из байткода технически невозможно реализовать.

Перечисленные проблемы делают необходимой разработку новой виртуальной машины, реализованной ближе к существующей среде исполнения языка. Такая виртуальная машина должна использовать представления значений Lama, взаимодействовать со средой исполнения, вызывать функции и поддерживать загрузку нескольких байткодных модулей. В рамках данной работы требуется доработать байткодное представление и разработать виртуальную машину для языка Lama, реализовать основные компоненты загрузки, связывания, валидации и исполнения, а также подготовить инфраструктуру проверки корректности полученного решения. Предлагаемый подход должен повысить полноту поддержки байткодного исполнения по сравнению с существующим интерпретатором, упростить дальнейшее развитие среды выполнения и создать основу для последующих исследований оптимизации исполнения программ.
